\documentclass[12pt,letterpaper]{article}
\usepackage{graphicx,textcomp}
\usepackage{natbib}
\usepackage{setspace}
\usepackage{fullpage}
\usepackage{color}
\usepackage[reqno]{amsmath}
\usepackage{amsthm}
\usepackage{fancyvrb}
\usepackage{amssymb,enumerate}
\usepackage[all]{xy}
\usepackage{endnotes}
\usepackage{lscape}
\newtheorem{com}{Comment}
\usepackage{float}
\usepackage{hyperref}
\newtheorem{lem} {Lemma}
\newtheorem{prop}{Proposition}
\newtheorem{thm}{Theorem}
\newtheorem{defn}{Definition}
\newtheorem{cor}{Corollary}
\newtheorem{obs}{Observation}
\usepackage[compact]{titlesec}
\usepackage{dcolumn}
\usepackage{tikz}
\usetikzlibrary{arrows}
\usepackage{multirow}
\usepackage{xcolor}
\newcolumntype{.}{D{.}{.}{-1}}
\newcolumntype{d}[1]{D{.}{.}{#1}}
\definecolor{light-gray}{gray}{0.65}
\usepackage{url}
\usepackage{listings}
\usepackage{color}

\definecolor{codegreen}{rgb}{0,0.6,0}
\definecolor{codegray}{rgb}{0.5,0.5,0.5}
\definecolor{codepurple}{rgb}{0.58,0,0.82}
\definecolor{backcolour}{rgb}{0.95,0.95,0.92}

\lstdefinestyle{mystyle}{
	backgroundcolor=\color{backcolour},   
	commentstyle=\color{codegreen},
	keywordstyle=\color{magenta},
	numberstyle=\tiny\color{codegray},
	stringstyle=\color{codepurple},
	basicstyle=\footnotesize,
	breakatwhitespace=false,         
	breaklines=true,                 
	captionpos=b,                    
	keepspaces=true,                 
	numbers=left,                    
	numbersep=5pt,                  
	showspaces=false,                
	showstringspaces=false,
	showtabs=false,                  
	tabsize=2
}
\lstset{style=mystyle}
\newcommand{\Sref}[1]{Section~\ref{#1}}
\newtheorem{hyp}{Hypothesis}

\title{Problem Set 2}
\date{Due: February 18, 2026}
\author{Applied Stats II}


\begin{document}
	\maketitle
	\section*{Instructions}
	\begin{itemize}
		\item Please show your work! You may lose points by simply writing in the answer. If the problem requires you to execute commands in \texttt{R}, please include the code you used to get your answers. Please also include the \texttt{.R} file that contains your code. If you are not sure if work needs to be shown for a particular problem, please ask.
		\item Your homework should be submitted electronically on GitHub in \texttt{.pdf} form.
		\item This problem set is due before 23:59 on Wednesday February 18, 2026. No late assignments will be accepted.

	\end{itemize}

	\vspace{.25cm}

\noindent We're interested in what types of international environmental agreements or policies people support (\href{https://www.pnas.org/content/110/34/13763}{Bechtel and Scheve 2013)}. So, we asked 8,500 individuals whether they support a given policy, and for each participant, we vary the (1) number of countries that participate in the international agreement and (2) sanctions for not following the agreement. \\

\noindent Load in the data labeled \texttt{climateSupport.RData} on GitHub, which contains an observational study of 8,500 observations.

\begin{itemize}
	\item
	Response variable: 
	\begin{itemize}
		\item \texttt{choice}: 1 if the individual agreed with the policy; 0 if the individual did not support the policy
	\end{itemize}
	\item
	Explanatory variables: 
	\begin{itemize}
		\item
		\texttt{countries}: Number of participating countries [20 of 192; 80 of 192; 160 of 192]
		\item
		\texttt{sanctions}: Sanctions for missing emission reduction targets [None, 5\%, 15\%, and 20\% of the monthly household costs given 2\% GDP growth]
		
	\end{itemize}
	
\end{itemize}

\newpage
\noindent Please answer the following questions:

\begin{enumerate}
	\item
	Remember, we are interested in predicting the likelihood of an individual supporting a policy based on the number of countries participating and the possible sanctions for non-compliance.
	\begin{enumerate}
		\item [] Fit an additive model. Provide the summary output, the global null hypothesis, and $p$-value. Please describe the results and provide a conclusion.

	\end{enumerate}
	
\texttt{RESPONSE}:\\

I have fitted an additive logistic regression model predicting individual support for a policy based on the number of participating countries and the level of sanctions:

\vspace{0.2in}
\lstinputlisting[language=R, firstline=71,lastline=73]{PS02_answers_RLP.R}
\vspace{0.2in}

The model's summary output is shown in \texttt{Table 1}.

\begin{table}[!htbp] \centering 
	\caption{} 
	\label{} 
	\begin{tabular}{@{\extracolsep{5pt}}lc} 
		\\[-1.8ex]\hline 
		\hline \\[-1.8ex] 
		& \multicolumn{1}{c}{\textit{Dependent variable:}} \\ 
		\cline{2-2} 
		\\[-1.8ex] & choice \\ 
		\hline \\[-1.8ex] 
		countries.L & 0.458$^{***}$ \\ 
		& (0.038) \\ 
		& \\ 
		countries.Q & $-$0.010 \\ 
		& (0.038) \\ 
		& \\ 
		sanctions.L & $-$0.276$^{***}$ \\ 
		& (0.044) \\ 
		& \\ 
		sanctions.Q & $-$0.181$^{***}$ \\ 
		& (0.044) \\ 
		& \\ 
		sanctions.C & 0.150$^{***}$ \\ 
		& (0.044) \\ 
		& \\ 
		Constant & $-$0.006 \\ 
		& (0.022) \\ 
		& \\ 
		\hline \\[-1.8ex] 
		Observations & 8,500 \\ 
		Log Likelihood & $-$5,784.130 \\ 
		Akaike Inf. Crit. & 11,580.260 \\ 
		\hline 
		\hline \\[-1.8ex] 
		\textit{Note:}  & \multicolumn{1}{r}{$^{*}$p$<$0.1; $^{**}$p$<$0.05; $^{***}$p$<$0.01} \\ 
	\end{tabular} 
\end{table} 

The results show a statistically significant positive effect of the number of participating countries on the outcome variable. This indicates a linear relationship between the number of countries involved and policy support: the more countries participate in the agreement, the higher the likelihood that individuals support it.\

As for sanctions, all three coefficients (linear, quadratic, and cubic) are statistically significant. Thus, \textit{ceteris paribus}, on average, higher levels of sanctions are associated with lower support for the policy, but this relationship is not strictly linear. Instead, the presence of higher-order terms suggests a non-linear pattern, and there may even be limited ranges in which support slightly increases, as indicated by the cubic component.

Finally, we run a likelihood ratio test (ANOVA) to assess whether we should reject the global null hypothesis and to obtain a p-value for the overall model:

\vspace{0.2in}
\lstinputlisting[language=R, firstline=61,lastline=64]{PS02_RLP.R}
\vspace{0.2in}

\begin{table}[!htbp] \centering 
	\caption{} 
	\label{} 
	\begin{tabular}{@{\extracolsep{5pt}}lccccc} 
		\\[-1.8ex]\hline 
		\hline \\[-1.8ex] 
		Statistic & \multicolumn{1}{c}{N} & \multicolumn{1}{c}{Mean} & \multicolumn{1}{c}{St. Dev.} & \multicolumn{1}{c}{Min} & \multicolumn{1}{c}{Max} \\ 
		\hline \\[-1.8ex] 
		Resid. Df & 2 & 8,496.500 & 3.536 & 8,494 & 8,499 \\ 
		Resid. Dev & 2 & 11,675.830 & 152.134 & 11,568.260 & 11,783.410 \\ 
		Df & 1 & 5.000 &  & 5 & 5 \\ 
		Deviance & 1 & 215.150 &  & 215.150 & 215.150 \\ 
		Pr(\textgreater Chi) & 1 & 0.000 &  & 0 & 0 \\ 
		\hline \\[-1.8ex] 
	\end{tabular} 
\end{table}

The results, as shown in \texttt{Table 2}, indicate that the global null hypothesis--that all coefficients are equal to zero, implying that neither the number of participating countries nor the level of sanctions affects support for the policy--can be rejected. The associated p-value is smaller than 0.001, indicating that the model is highly statistically significant overall.
	
	\item
	If any of the explanatory variables are statistically significant in this model, then:
	\begin{enumerate}
		\item
		For the policy in which nearly all countries participate [160 of 192], how does increasing sanctions from 5\% to 15\% change the odds that an individual will support the policy? (Interpretation of a coefficient)\\

\texttt{RESPONSE}:\\

\begin{lstlisting}[language=R]
	> p1_q2a_160 <- predict(model_q1,
	      newdata = data.frame(countries="160 of 192", sanctions="5%"), type="link")
	> p2_q2a_160 <- predict(model_q1,
	             newdata = data.frame(countries="160 of 192", sanctions="15%"), type="link")
	> answer_q2_a <- exp(p2_q2a_160 - p1_q2a_160)
	> print(answer_q2_a)
	1 
	0.7224531 \end{lstlisting}
\vspace{0.2in}

When almost all countries participate (160 out of 192), increasing sanctions from 5\% to 15\% multiplies the odds of supporting the policy by 0.72, meaning the odds decrease by about 28\% and individuals are less likely to support it.\\

		\item
		For the policy in which very few countries participate [20 of 192], how does increasing sanctions from 5\% to 15\% change the odds that an individual will support the policy? (Interpretation of a coefficient)\\
		
\texttt{RESPONSE}:\\

\begin{lstlisting}[language=R]
	> p1_q2_20 <- predict(model_q1, newdata = data.frame(countries="20 of 192", sanctions="5%"), type="link")
	> p2_q2_20 <- predict(model_q1, newdata = data.frame(countries="20 of 192", sanctions="15%"), type="link")
	> answer_q2_b <- exp(p2_q2_20 - p1_q2_20)
	> print(answer_q2_b)
	1 
	0.7224531  \end{lstlisting}
\vspace{0.2in}

When only a reduced number of countries participates (20 out of 192), increasing sanctions from 5\% to 15\% again multiplies the odds of supporting the policy by 0.72, meaning the odds decrease by about 28\% and individuals are less likely to support it.\\	

		\item
		What is the estimated probability that an individual will support a policy if there are 80 of 192 countries participating with no sanctions?\\
		
\texttt{RESPONSE}:\\

\begin{lstlisting}[language=R]
	> p1_q3_80 <- predict(model_q1, newdata = data.frame(countries="80 of 192", sanctions="None"), type="response")
	> print(p1_q3_80)
	1 
	0.5159191  \end{lstlisting}
\vspace{0.2in}
		
The estimated probability that an individual supports the policy when 80 out of 192 countries participate and there are no sanctions is around a half (about 51.6\%).\\
	
	\end{enumerate}
	\item
	Would the answers to 2a and 2b potentially change if we included an interaction term in this model? Why?\\
	
\texttt{RESPONSE}:\\
	
	Right now, the model is additive, assuming that the effect of "sanctions" on "choice" is the same for all levels of "countries". However, it is likely that this effect varies depending on how many countries are involved in the policy. If we use an interactive model, we might see that the effect of one variable (in this case, "sanctions") depends on the level of another ("countries").\\
	
	Therefore, including an interaction term may change the answers to 2a and 2b.
	
	\begin{itemize}
		\item Perform a test to see if including an interaction is appropriate.\\
		
\texttt{RESPONSE}:\\

To perform this test we first need to generate an interaction model that creates an interaction term between "countries" and "sanctions":\\

\begin{lstlisting}[language=R]
	model_q3 <- glm(choice ~ countries * sanctions, data=climateSupport, family=binomial)
 \end{lstlisting}
\vspace{0.2in}

Then, we must run an ANOVA test to see if this model is more appropriate compared to the one we already have (additive):

\begin{lstlisting}[language=R]
	anova_test_q3 <- anova(model_q1, model_q3, test = "Chisq")
\end{lstlisting}
\vspace{0.2in}

The results of the test are shown in \texttt{Table 3}.

\begin{table}[!htbp] \centering 
	\caption{} 
	\label{} 
	\begin{tabular}{@{\extracolsep{5pt}}lccccc} 
		\\[-1.8ex]\hline 
		\hline \\[-1.8ex] 
		Statistic & \multicolumn{1}{c}{N} & \multicolumn{1}{c}{Mean} & \multicolumn{1}{c}{St. Dev.} & \multicolumn{1}{c}{Min} & \multicolumn{1}{c}{Max} \\ 
		\hline \\[-1.8ex] 
		Resid. Df & 2 & 8,491.000 & 4.243 & 8,488 & 8,494 \\ 
		Resid. Dev & 2 & 11,565.110 & 4.450 & 11,561.970 & 11,568.260 \\ 
		Df & 1 & 6.000 &  & 6 & 6 \\ 
		Deviance & 1 & 6.293 &  & 6.293 & 6.293 \\ 
		Pr(\textgreater Chi) & 1 & 0.391 &  & 0.391 & 0.391 \\ 
		\hline \\[-1.8ex] 
	\end{tabular} 
\end{table} 

The ANOVA test shows a p-value of 0.39, well above the standards of statistical significance. In other words, the interaction between "countries" and "sanctions" is not statistically significant. This means that adding an interaction term does not really improve the model fit. Therefore, the original additive model is sufficient.

	\end{itemize}
	\end{enumerate}


\end{document}
